\section{FunQA Extension Datasets}\label{appendix_b: FunQA Extended dataset}
Our main objective in designing these extended datasets is to leverage our high-quality annotated data and provide a rich and suitable data format for models. In addition, we also wanted to test the capability of GPT-3.5 and the quality of the dataset, and it turns out that GPT-3.5 can expand our data tens of times after providing high-quality free-text annotations, and our FunQA dataset is extremely scalable.
\begin{figure}[htbp]
    \centering
    \includegraphics[width=\linewidth]{figures_appendix/prompt-mc.pdf}
    \vspace{-0.6cm}
    \caption{\textbf{GPT prompt for FunQA Multi-choice.} We give GPT-3.5 prompt as system, and then enter the description and explanation of the video, and GPT-3.5 will output the data in the form of multiple-choice questions according to the prompt.}
    \vspace{-0.4cm}
    \label{fig:prompt-mcqa}
\end{figure}

\begin{figure}[htbp]
    \centering
    \includegraphics[width=\linewidth]{figures_appendix/prompt-dia.pdf}
    \vspace{-0.6cm}
    \caption{\textbf{GPT prompt for FunQA Dialog.} We give GPT prompt as system, and then input the description and explanation of the video, and GPT-3.5 will output the data in the form of dialogue according to the prompt.}
    \vspace{-0.4cm}
    \label{fig:prompt-dia}
\end{figure}
\subsection{FunQA Multi-choice Dataset} \label{appendix_b: FunQA MC}
FunQA Multi-choice Dataset is prepared to provide training and testing for arbitrary models, in this dataset our QA pairs are in the form of multiple choice, the answer is a word, phrase, or short sentence, and the type of questions are all descriptions. An example of the prompt given to GPT-3.5 and the data generated can be seen in Figure \ref{fig:prompt-mcqa}.
\subsection{FunQA Dialog Dataset}\label{appendix_b: FunQA Dia}
Most of the current LLMs are in the form of dialogues. To cater to their data input, we produced the FunQA Dialog dataset, in which we used GPT-3.5 to convert QA pairs into recursive dialogues with added context. In addition to this, we also let GPT-3.5 freely generate future dialogue development, which greatly expands the amount of information. An example of the prompt given to GPT-3.5 and the data generated can be seen in Figure \ref{fig:prompt-dia}.
